% ============================================================================
%  Atlas - POC Helpdesk GLPI / PRA-DR : Procedure d'exploitation
%  Compilation : pdflatex POC_Atlas_Procedure.tex  (2 passes pour la TOC)
%  Paquets requis (TeX Live full / MiKTeX) : babel, geometry, xcolor,
%  listings, tcolorbox, booktabs, tabularx, titlesec, fancyhdr, hyperref.
% ============================================================================
\documentclass[11pt,a4paper]{article}

\usepackage[T1]{fontenc}
\usepackage[utf8]{inputenc}
\usepackage{lmodern}
\usepackage[french]{babel}
\usepackage[a4paper,margin=2.4cm]{geometry}
\usepackage{xcolor}
\usepackage{graphicx}
\usepackage{booktabs}
\usepackage{tabularx}
\usepackage{array}
\usepackage{listings}
\usepackage[most]{tcolorbox}
\usepackage{titlesec}
\usepackage{fancyhdr}
\usepackage{enumitem}
\usepackage{underscore} % permet d'ecrire les _ des noms de fichiers en clair
\usepackage[hidelinks,colorlinks=true,linkcolor=atlasblue,urlcolor=atlasblue]{hyperref}

% ---- Palette Atlas (NOC / mission control) ----
\definecolor{atlasblue}{HTML}{1F6FEB}
\definecolor{atlasdark}{HTML}{0D1B2A}
\definecolor{atlasgray}{HTML}{F2F4F8}
\definecolor{atlasgreen}{HTML}{2EA043}
\definecolor{atlasred}{HTML}{CF222E}
\definecolor{atlasamber}{HTML}{9A6700}
\definecolor{codebg}{HTML}{F6F8FA}
\definecolor{codekw}{HTML}{CF222E}
\definecolor{codecom}{HTML}{6E7781}
\definecolor{codestr}{HTML}{0A3069}

% ---- Titres colores ----
\titleformat{\section}{\Large\bfseries\color{atlasblue}}{\thesection.}{0.6em}{}
\titleformat{\subsection}{\large\bfseries\color{atlasdark}}{\thesubsection}{0.6em}{}
\titlespacing*{\section}{0pt}{2.2ex plus 1ex minus .2ex}{1.2ex}

% ---- En-tetes / pieds ----
\pagestyle{fancy}
\fancyhf{}
\lhead{\small\color{atlasdark}\textbf{Atlas} — POC Helpdesk / PRA-DR}
\rhead{\small\color{codecom}\thepage}
\rfoot{\small\color{codecom}Document interne — usage POC}
\renewcommand{\headrulewidth}{0.4pt}

% ---- Style code ----
\lstset{
  basicstyle=\ttfamily\small,
  backgroundcolor=\color{codebg},
  frame=leftline,
  framerule=2pt,
  rulecolor=\color{atlasblue},
  keywordstyle=\color{codekw}\bfseries,
  commentstyle=\color{codecom}\itshape,
  stringstyle=\color{codestr},
  showstringspaces=false,
  breaklines=true,
  breakatwhitespace=true,
  columns=fullflexible,
  keepspaces=true,
  xleftmargin=1em,
  aboveskip=1em,
  belowskip=1em,
  inputencoding=utf8,
  extendedchars=true,
  literate=%
    {é}{{\'e}}1 {è}{{\`e}}1 {ê}{{\^e}}1 {ë}{{\"e}}1
    {à}{{\`a}}1 {â}{{\^a}}1 {ä}{{\"a}}1
    {î}{{\^i}}1 {ï}{{\"i}}1 {ô}{{\^o}}1 {ö}{{\"o}}1
    {û}{{\^u}}1 {ù}{{\`u}}1 {ü}{{\"u}}1 {ç}{{\c c}}1
    {É}{{\'E}}1 {È}{{\`E}}1 {Ê}{{\^E}}1 {À}{{\`A}}1
    {«}{{\guillemotleft}}1 {»}{{\guillemotright}}1
    {→}{{$\rightarrow$}}1 {≤}{{$\leq$}}1 {œ}{{\oe}}1
}
\lstdefinestyle{bash}{language=bash,morekeywords={ansible,ansible-playbook,ssh,rsync,systemctl,mysql,mysqldump,zcat,pct,iptables,curl}}
\lstdefinestyle{yaml}{language=yaml}
\lstdefinestyle{plain}{}

% ---- Boites d'appel ----
\newtcolorbox{pourquoi}{enhanced,breakable,colback=atlasblue!5,colframe=atlasblue,
  fonttitle=\bfseries,title=Pourquoi cette commande ?,left=2mm,right=2mm,top=1mm,bottom=1mm}
\newtcolorbox{attention}{enhanced,breakable,colback=atlasred!5,colframe=atlasred,
  fonttitle=\bfseries,title=Attention,left=2mm,right=2mm,top=1mm,bottom=1mm}
\newtcolorbox{astuce}{enhanced,breakable,colback=atlasgreen!6,colframe=atlasgreen,
  fonttitle=\bfseries,title=Astuce,left=2mm,right=2mm,top=1mm,bottom=1mm}

\newcommand{\cmd}[1]{{\ttfamily #1}}

\begin{document}

% ======================= PAGE DE TITRE =======================
\begin{titlepage}
\centering
\vspace*{2cm}
{\color{atlasblue}\rule{\linewidth}{2pt}}\\[0.6cm]
{\Huge\bfseries\color{atlasdark} Projet Atlas}\\[0.35cm]
{\LARGE Portail Helpdesk GLPI — Plan de Reprise d'Activité}\\[0.25cm]
{\large\color{codecom} Procédure d'exploitation du POC (Infrastructure as Code)}\\[0.4cm]
{\color{atlasblue}\rule{\linewidth}{2pt}}\\[1.4cm]

\begin{tcolorbox}[width=0.85\linewidth,colback=atlasgray,colframe=atlasdark,sharp corners]
\centering\small
\textbf{Contraintes de reprise} : RPO $\leq$ 20 min \quad|\quad RTO $\leq$ 40 min\\[2pt]
\textbf{Méthode} : 100\% IaC (Ansible) — aucun « clic-clic » non tracé\\[2pt]
\textbf{Hyperviseur} : Proxmox VE 8.4 (pve01, 138.201.135.108)\\[2pt]
\textbf{Périmètre} : 9 conteneurs LXC — 3 VLAN (DMZ / LAN-ADMIN / MGMT)
\end{tcolorbox}

\vfill
{\color{codecom}\small Document interne — support de présentation et d'exploitation du POC}\\
{\color{codecom}\small Généré le \today}
\end{titlepage}

\tableofcontents
\newpage

% ======================= 1. PRESENTATION =======================
\section{Présentation du POC}

\subsection{Objectif}
Atlas est un \textbf{portail de tickets d'assistance (GLPI)} déployé de bout en bout par
Infrastructure as Code. L'objectif du POC est de démontrer qu'un incident majeur
(perte de VM, corruption de données, panne de service) peut être détecté puis corrigé
\textbf{dans les délais contractuels} : RPO $\leq$ 20 minutes (perte de données maximale)
et RTO $\leq$ 40 minutes (durée de remise en service).

\subsection{Architecture}
Deux chemins réseau \emph{indépendants} coexistent :

\begin{itemize}[leftmargin=1.4em]
  \item \textbf{Chemin utilisateur (web)} : Internet $\rightarrow$ Reverse Proxy TLS
        (atlrp01p, DMZ) $\rightarrow$ GLPI (atlapp01p) $\rightarrow$ MariaDB (atldb01p).
  \item \textbf{Chemin administration (SSH)} : Admin $\rightarrow$ Bastion (atlbst01p, MGMT)
        $\rightarrow$ rebond vers tous les LXC. Aucun service n'est exposé sur Internet ;
        l'accès admin passe par tunnels SSH.
\end{itemize}

\begin{table}[h]
\centering\small
\renewcommand{\arraystretch}{1.2}
\begin{tabularx}{\linewidth}{@{}l l l X@{}}
\toprule
\textbf{Conteneur} & \textbf{VMID} & \textbf{IP} & \textbf{Rôle} \\
\midrule
atlbst01p & 200 & 10.90.0.10 & Bastion SSH — point d'entrée admin unique \\
atlrp01p  & 201 & 10.20.0.10 & Reverse proxy Nginx (TLS) \\
atlapp01p & 202 & 10.20.0.11 & Application GLPI (PHP-FPM) \\
atldb01p  & 203 & 10.30.0.10 & Base de données MariaDB \\
atlsmtp01p& 204 & 10.30.0.11 & Relais SMTP Postfix \\
atlzbx01p & 205 & 10.30.0.12 & Supervision Zabbix (serveur + frontend) \\
atlbkp01p & 206 & 10.30.0.13 & Serveur de sauvegarde (copies distantes) \\
atlgrf01p & 207 & 10.30.0.14 & Grafana (visualisation) \\
atltst01p & 208 & 10.20.0.20 & Machine de test (chaos engineering) \\
\bottomrule
\end{tabularx}
\caption{Les 9 conteneurs LXC du POC Atlas.}
\end{table}

\begin{table}[h]
\centering\small
\begin{tabularx}{\linewidth}{@{}l l X@{}}
\toprule
\textbf{VLAN} & \textbf{Sous-réseau} & \textbf{Usage} \\
\midrule
20 (DMZ)        & 10.20.0.0/24 & Services exposés indirectement (reverse proxy, GLPI, test) \\
30 (LAN-ADMIN)  & 10.30.0.0/24 & Données \& supervision (BDD, SMTP, Zabbix, backup, Grafana) \\
90 (MGMT)       & 10.90.0.0/24 & Management — bastion SSH uniquement \\
\bottomrule
\end{tabularx}
\caption{Segmentation réseau. Le transit inter-VLAN est refusé par défaut (politique iptables \cmd{FORWARD DROP}), seuls les flux de la matrice de moindre privilège sont autorisés.}
\end{table}

% ======================= 2. PREREQUIS =======================
\section{Prérequis et poste de pilotage}

Toutes les commandes Ansible se lancent depuis le \textbf{control node}, l'hyperviseur
\textbf{pve01} (\cmd{138.201.135.108}), où réside le dépôt du projet.

\begin{lstlisting}[style=bash]
# Se connecter au control node
ssh root@138.201.135.108
cd /root/atlas          # repertoire du depot Ansible

# Verifier les outils
ansible --version        # Ansible core
ansible-inventory -i develop/hosts.yaml --graph   # voir les groupes/hotes
\end{lstlisting}

\begin{pourquoi}
La commande \cmd{ansible-inventory --graph} affiche l'arborescence des groupes
(\cmd{proxmox}, \cmd{mgmt}, \cmd{dmz}, \cmd{lan_admin}) et confirme que l'inventaire
est lisible \emph{avant} de lancer quoi que ce soit. C'est le premier filet de sécurité :
si un hôte manque ici, le playbook échouera.
\end{pourquoi}

\paragraph{Éléments sensibles (hors dépôt Git).}
\begin{itemize}[leftmargin=1.4em]
  \item Clé SSH d'administration : \cmd{/root/.ssh/atlas_ed25519} (rebond bastion).
  \item Mot de passe du coffre Ansible : \cmd{.vault_pass} (déchiffre les secrets \cmd{vault\_*}).
  \item Configuration : \cmd{ansible.cfg} pointe déjà l'inventaire, le \cmd{roles_path}
        et le vault. Aucune option à répéter en ligne de commande en usage normal.
\end{itemize}

\begin{attention}
\textbf{Reprise du control node.} Si \cmd{pve01} est lui-même le sinistré, Ansible doit
pouvoir tourner depuis un poste de secours : copier \cmd{atlas_ed25519} + \cmd{.vault_pass}
(coffre hors site) puis surcharger la clé sans éditer \cmd{ansible.cfg} :
\cmd{ansible-playbook ... -e ansible_ssh_private_key_file=<chemin>}. Cette bascule doit
être testée régulièrement pour garantir le RTO.
\end{attention}

% ======================= 3. FIX RECENTS =======================
\section{Corrections intégrées à l'IaC (juillet 2026)}
Trois corrections de supervision ont été appliquées puis \textbf{pérennisées dans le
playbook} (elles survivent donc à toute réexécution) :

\begin{table}[h]
\centering\small
\begin{tabularx}{\linewidth}{@{}l X l@{}}
\toprule
\textbf{Fix} & \textbf{Correction} & \textbf{Fichier modifié} \\
\midrule
B & Flux firewall supervision $\rightarrow$ bastion (VLAN30$\rightarrow$90, ports 10050/10051)
  & \cmd{roles/proxmox_network/defaults/main.yaml} \\
C & Agent Zabbix accepte les checks passifs locaux (\cmd{127.0.0.1})
  & \cmd{roles/zabbix_agent/templates/zabbix_agent2.conf.j2} \\
A & Désactivation des items internes « not supported » (bruit)
  & \cmd{roles/zabbix/tasks/tune_health.yaml} \\
\bottomrule
\end{tabularx}
\caption{Corrections encodées dans l'IaC. Un simple \cmd{--tags proxmox\_network,zabbix} redéploie l'ensemble.}
\end{table}

% ======================= 4. RELANCER LE PLAYBOOK EN SECURITE =======================
\section{Relancer le playbook — la séquence de tests \emph{avant} d'appliquer}
\label{sec:tests}

\begin{astuce}
Règle d'or : on ne lance jamais un \cmd{ansible-playbook} « à sec » en production sans
avoir passé les 3 filets ci-dessous, dans l'ordre : \textbf{(1)} syntaxe,
\textbf{(2)} périmètre, \textbf{(3)} simulation. Puis seulement l'application réelle.
\end{astuce}

\subsection{Étape 1 — Vérification de syntaxe}
\begin{lstlisting}[style=bash]
ansible-playbook site.yaml --syntax-check
\end{lstlisting}
\begin{pourquoi}
Analyse le YAML et le graphe des rôles \emph{sans se connecter} à aucune machine.
Détecte les fautes d'indentation, les tâches malformées, les rôles introuvables.
Coût nul, à lancer après chaque édition du playbook.
\end{pourquoi}

\subsection{Étape 2 — Vérifier le périmètre (quoi et où)}
\begin{lstlisting}[style=bash]
ansible-playbook site.yaml --list-hosts     # sur quels hotes ca va tourner
ansible-playbook site.yaml --list-tasks     # quelles taches, dans quel ordre
\end{lstlisting}
\begin{pourquoi}
\cmd{--list-hosts} garantit qu'on ne va pas toucher un hôte non prévu ;
\cmd{--list-tasks} montre l'enchaînement exact (utile pour repérer où agit un tag).
Toujours confirmer la cible \emph{avant} d'exécuter, surtout en production.
\end{pourquoi}

\subsection{Étape 3 — Simulation (dry-run) avec diff}
\begin{lstlisting}[style=bash]
# Simulation globale : rien n'est modifie, on voit ce QUI changerait
ansible-playbook site.yaml --check --diff

# Simulation ciblee sur un seul role + un seul hote
ansible-playbook site.yaml --check --diff --tags zabbix --limit atlzbx01p
\end{lstlisting}
\begin{pourquoi}
\cmd{--check} exécute en « mode blanc » (aucune écriture réelle) et \cmd{--diff} affiche
les différences ligne à ligne des fichiers qui seraient modifiés. C'est la répétition
générale : si le diff montre autre chose que prévu, on corrige \emph{avant} d'appliquer.
\end{pourquoi}
\begin{attention}
Le mode \cmd{--check} n'est pas parfait : les tâches qui dépendent d'un résultat
précédent (API, \cmd{shell}) peuvent être « skippées » ou signalées à tort. Il valide
l'intention, pas un résultat garanti à 100\%.
\end{attention}

\subsection{Étape 4 — Application réelle}
\begin{lstlisting}[style=bash]
# Deploiement complet (idempotent : rejouable sans risque)
ansible-playbook site.yaml

# Application ciblee (recommande pour une correction precise)
ansible-playbook site.yaml --tags zabbix                    # supervision seule
ansible-playbook site.yaml --tags proxmox_network           # firewall/VLAN de l'hote
ansible-playbook site.yaml --tags zabbix --limit atlzbx01p  # 1 role + 1 hote
\end{lstlisting}
\begin{pourquoi}
Le playbook est \textbf{idempotent} : le relancer sur une infra déjà conforme ne change
rien (état « ok », zéro « changed »). Les \textbf{tags} permettent de n'appliquer qu'un
sous-ensemble (ex. seulement la supervision), et \cmd{--limit} restreint aux hôtes visés :
c'est le principe de moindre impact.
\end{pourquoi}

\begin{table}[h]
\centering\small
\begin{tabularx}{\linewidth}{@{}l X@{}}
\toprule
\textbf{Tag} & \textbf{Périmètre} \\
\midrule
\cmd{proxmox_network}   & Réseau VLAN + firewall iptables de l'hyperviseur \\
\cmd{proxmox_provision} & Création/configuration des conteneurs LXC \\
\cmd{common}            & Socle commun (durcissement SSH, UFW, NTP) \\
\cmd{bastion}           & Bastion SSH \\
\cmd{mariadb}           & Base de données \\
\cmd{backup}            & Chaîne de sauvegarde \\
\cmd{glpi} / \cmd{glpi_seed} & Application GLPI / données de démo \\
\cmd{reverse_proxy}     & Reverse proxy Nginx TLS \\
\cmd{zabbix}            & Serveur + agents de supervision \\
\cmd{grafana}           & Tableaux de bord \\
\cmd{testlab}           & Machine de chaos engineering \\
\bottomrule
\end{tabularx}
\caption{Tags disponibles dans \cmd{site.yaml} (voir l'en-tête du fichier).}
\end{table}

% ======================= 5. PRA RESTORE FULL =======================
\section{PRA — Restauration complète (perte de VM)}
Scénario : un ou plusieurs conteneurs sont perdus. Le playbook \cmd{pra_restore_full.yaml}
reconstruit toute la chaîne \emph{dans l'ordre de dépendance}, restaure les données, puis
\textbf{mesure le RTO} automatiquement.

\begin{lstlisting}[style=bash]
# 1) Tests prealables (comme en section 4)
ansible-playbook pra_restore_full.yaml --syntax-check
ansible-playbook pra_restore_full.yaml --list-tasks

# 2) Restauration complete (dernier dump distant utilise automatiquement)
ansible-playbook pra_restore_full.yaml

# 3) Restauration a partir d'un dump precis (optionnel)
ansible-playbook pra_restore_full.yaml \
  -e "pra_restore_dump=/var/backups/atlas/db-remote/glpi_dump_20260715_195001.sql.gz"
\end{lstlisting}

\begin{pourquoi}
Sans \cmd{pra_restore_dump}, le playbook sélectionne \emph{automatiquement} le dump le plus
récent présent sur \cmd{atlbkp01p} (\cmd{ls -t | head -1}) : on repart donc du point de
récupération le plus proche possible (RPO minimal). On ne fournit un dump précis que pour
rejouer un état antérieur particulier.
\end{pourquoi}

\paragraph{Ce que fait le playbook (résumé).}
\begin{enumerate}[leftmargin=1.6em]
  \item Horodatage de début + journal \cmd{/var/log/atlas/pra_restore_full.log}.
  \item Reprovisioning LXC, socle \cmd{common}, bastion.
  \item Reconstruction des services : MariaDB $\rightarrow$ SMTP $\rightarrow$ GLPI
        $\rightarrow$ reverse proxy $\rightarrow$ Zabbix $\rightarrow$ backup.
  \item Restauration des données (bloc \cmd{block/rescue/always}) : transfert du dump
        \cmd{atlbkp01p} $\rightarrow$ \cmd{atldb01p}, import dans \cmd{glpi_db},
        restauration du filestore GLPI, puis \textbf{suppression du dump temporaire}
        (jamais de base en clair dans \cmd{/tmp}).
  \item Vérifications : GLPI en HTTP direct \emph{et} via le reverse proxy HTTPS,
        comptage des tickets, port Zabbix 10051 ouvert.
  \item Calcul et journalisation du \textbf{RTO mesuré} (comparé à la cible 40 min).
\end{enumerate}

\begin{attention}
En cas d'échec pendant la restauration des données, le bloc \cmd{rescue} journalise
l'\textbf{ÉCHEC} et \cmd{any_errors_fatal} interrompt le PRA : on ne poursuit jamais sur
une base à moitié restaurée. Le \cmd{always} nettoie systématiquement le dump temporaire.
\end{attention}

% ======================= 6. PRA RESTORE GRANULAIRE =======================
\section{PRA — Restauration granulaire (tickets supprimés)}
Scénario : des tickets ont été supprimés/corrompus, mais l'infrastructure est saine.
On ne veut restaurer \emph{que} les tickets manquants, \textbf{sans écraser} ceux créés
depuis l'incident.

\begin{lstlisting}[style=bash]
# 1) Choisir le dump anterieur a l'incident (sur atldb01p)
ssh root@10.30.0.10 'ls -lht /var/backups/atlas/db/ | head'

# 2) Lancer la restauration granulaire (dump OBLIGATOIRE)
ansible-playbook pra_restore_granular.yaml \
  -e "pra_target_dump=/var/backups/atlas/db/glpi_dump_20260715_193001.sql.gz"
\end{lstlisting}

\begin{pourquoi}
La variable \cmd{pra_target_dump} est \textbf{obligatoire} (le playbook s'arrête via un
\cmd{assert} si elle manque) : restaurer des tickets est une opération chirurgicale, on
doit choisir explicitement la photo à réinjecter. Aucun automatisme ici, contrairement au
restore complet.
\end{pourquoi}

\paragraph{Mécanisme de sûreté.} Le dump est restauré dans une base \emph{temporaire}
(\cmd{glpi_restore_tmp}), jamais en production directement. Seules les tables de tickets
sont exportées avec \cmd{mysqldump --no-create-info --insert-ignore}, puis fusionnées :
les \cmd{INSERT IGNORE} ne touchent pas les lignes déjà présentes. \textbf{Les tickets
créés après l'incident ne sont donc jamais écrasés.} La base temporaire est supprimée en fin
de course, et le RPO effectif (âge du dump) est journalisé.

% ======================= 7. SAUVEGARDES =======================
\section{Sauvegardes et RPO}
Les dumps GLPI sont produits \textbf{toutes les 15 minutes} (cron
\cmd{5,20,35,50 * * * *}), ce qui garantit un RPO $\leq$ 15 min (meilleur que la cible de
20 min). Chaque dump est vérifié (\cmd{gzip -t}) puis répliqué hors-site sur \cmd{atlbkp01p}.

\begin{lstlisting}[style=bash]
# Verifier les derniers dumps locaux (sur la BDD)
ssh root@10.30.0.10 'ls -lht /var/backups/atlas/db/ | head'

# Verifier la replication hors-site (sur le backup)
ssh root@10.30.0.13 'ls -lht /var/backups/atlas/db-remote/ | head'

# Consulter le journal de dump
ssh root@10.30.0.10 'tail -5 /var/log/atlas/dump_glpi.log'
\end{lstlisting}

\begin{pourquoi}
Vérifier que le dump le plus récent est bien présent \emph{des deux côtés} (local +
hors-site) et daté de moins de 15 minutes prouve que la chaîne de sauvegarde est vivante —
c'est la condition sine qua non pour tenir le RPO en cas de sinistre.
\end{pourquoi}

% ======================= 8. CHAOS =======================
\section{Chaos engineering (validation de la détection)}
La machine de lab \cmd{atltst01p} embarque \cmd{chaos.sh} : un générateur d'incidents
\textbf{contrôlés et réversibles}, détectés par Zabbix. Il sert à prouver que la supervision
et les alertes fonctionnent réellement.

\begin{lstlisting}[style=bash]
# Sur atltst01p (via rebond bastion) :
/opt/atlas/scripts/chaos.sh run      # un scenario ALEATOIRE (cpu | disk | service)
/opt/atlas/scripts/chaos.sh cpu      # saturation CPU temporaire (stress-ng)
/opt/atlas/scripts/chaos.sh disk     # remplissage disque controle
/opt/atlas/scripts/chaos.sh service  # arret de nginx
/opt/atlas/scripts/chaos.sh clean    # REPARATION : retour a l'etat nominal

# Depuis le control node, en une ligne (ad-hoc Ansible) :
ansible atltst01p -m command -a "/opt/atlas/scripts/chaos.sh service"
ansible atltst01p -m command -a "/opt/atlas/scripts/chaos.sh clean"
\end{lstlisting}

\begin{pourquoi}
Déclencher un incident connu permet de vérifier que Zabbix \emph{voit} la panne (alerte),
puis que \cmd{clean} rétablit l'état nominal. C'est la boucle « je casse $\rightarrow$ ça
alerte $\rightarrow$ je répare » qui valide la supervision de bout en bout.
\end{pourquoi}

\begin{attention}
\cmd{chaos.sh} contient un \textbf{garde-fou} : il refuse de s'exécuter ailleurs que sur
l'hôte de lab attendu (\cmd{hostname -s}). Sans cela, une erreur d'inventaire pourrait
arrêter Nginx sur le vrai reverse proxy. Un cron déclenche aussi un incident toutes les
2 h, réparé automatiquement 45 min plus tard.
\end{attention}

% ======================= 9. SUPERVISION / ACCES =======================
\section{Accès aux interfaces (tunnels SSH)}
Aucun service n'est exposé sur Internet : on y accède par \textbf{tunnel SSH} depuis son
poste. Ouvrir le tunnel, puis pointer le navigateur sur \cmd{localhost}.

\begin{lstlisting}[style=bash]
# GLPI (via reverse proxy TLS)
ssh -L 8443:10.20.0.10:443 root@138.201.135.108     # -> https://localhost:8443

# Zabbix (supervision)
ssh -L 8080:10.30.0.12:8080 root@138.201.135.108    # -> http://localhost:8080

# Grafana (tableaux de bord)
ssh -L 3000:10.30.0.14:3000 root@138.201.135.108    # -> http://localhost:3000
\end{lstlisting}

\begin{pourquoi}
Le tunnel \cmd{ssh -L} crée un pont chiffré entre un port de votre poste et le service
interne, sans jamais l'exposer publiquement. La surface d'attaque reste minimale : seul le
port SSH de l'hôte est joignable de l'extérieur.
\end{pourquoi}

% ======================= 10. RECAP =======================
\section{Aide-mémoire — commandes essentielles}
\begin{table}[h]
\centering\small
\renewcommand{\arraystretch}{1.25}
\begin{tabularx}{\linewidth}{@{}X l@{}}
\toprule
\textbf{Objectif} & \textbf{Commande} \\
\midrule
Vérifier la syntaxe        & \cmd{ansible-playbook site.yaml --syntax-check} \\
Voir hôtes / tâches        & \cmd{ansible-playbook site.yaml --list-hosts / --list-tasks} \\
Simuler (dry-run)          & \cmd{ansible-playbook site.yaml --check --diff} \\
Déployer tout              & \cmd{ansible-playbook site.yaml} \\
Déployer un rôle           & \cmd{ansible-playbook site.yaml --tags zabbix} \\
Restauration complète      & \cmd{ansible-playbook pra_restore_full.yaml} \\
Restauration granulaire    & \cmd{ansible-playbook pra_restore_granular.yaml -e "pra_target_dump=..."} \\
Vérifier les sauvegardes   & \cmd{ssh root@10.30.0.10 'ls -lht /var/backups/atlas/db/'} \\
Lancer un incident chaos   & \cmd{ansible atltst01p -m command -a ".../chaos.sh run"} \\
Réparer le chaos           & \cmd{ansible atltst01p -m command -a ".../chaos.sh clean"} \\
\bottomrule
\end{tabularx}
\caption{Récapitulatif opérationnel.}
\end{table}

\vfill
\begin{center}
{\color{atlasblue}\rule{0.5\linewidth}{1pt}}\\[4pt]
{\small\color{codecom} Atlas — POC Helpdesk / PRA-DR — Infrastructure as Code}
\end{center}

\end{document}
